% mass-gap-theorem.tex — Chapter on the SU(2) mass-gap theorem,
% composing Phases 7.1, 8.1, 8.2 into the Clay-statement form.
%
% Linked code / Lean:
%   - computations/mass_gap_bound.py (derive_mass_gap_bound)
%   - lean/YMGap/MassGap/SU2.lean    (mass_gap_SU2 headline theorem)
%   - \Cref{ch:mass-gap-theorem} of work_plan.md.

\chapter{The \texorpdfstring{$\SUtwo$}{SU(2)} Mass-Gap Theorem:
  A Constructive Lower Bound from
  \texorpdfstring{$\sighinf > 0$}{sigma\_infty > 0}}
\label{ch:mass-gap-theorem}

\noindent\emph{The headline theorem: $\sigma(\Hop) \cap (0, \massgap)
= \emptyset$, with $\massgap \in \Qp$ derived from the Phase-3 pilot
numerics and the OS reconstruction of Chapter~\ref{ch:os-reconstruction}.}

We state and prove the Osterwalder--Schrader-reconstruction mass-gap
theorem for $\SUtwo$ Yang--Mills theory.  The proof composes the
three preceding milestones into the Clay-statement form: the
$\Qp$-rational lower bound $\sighinf \ge \varepsilon > 0$ from
the $\sighinf$-positivity theorem, the verified OS4 cluster property with explicit
exponential rate from \Cref{ch:os-axioms}
(Chapter~\ref{ch:os-axioms}), and the OS reconstruction
quintuple with the cluster-to-gap bridge from \Cref{ch:os-reconstruction}
(Chapter~\ref{ch:os-reconstruction}), together imply that the
continuum-limit Hamiltonian $\Hop$ on $\Hphys$ satisfies
\[
  \sigma(\Hop) \cap (0, \massgap) \;=\; \emptyset
\]
for an explicit $\massgap > 0$ derived from the Phase-3 pilot
numerics.  The constructed $\massgap$ is a strict (not sharp) lower
bound on the mass gap, presented in two flavours: canonical
($g_0 = 0.30$ pilot, $\massgap \ge 999/1000$) and Wilson
($g_0 = 0.50$ pilot, $\massgap \ge 999/1000$), both in dimensionless
units where $\sigma_0^{1/2}$ sets the scale at lattice scale $0$.

\section{The Clay-statement mass-gap theorem for $\SUtwo$}
\label{sec:mass-gap-clay}

\subsection{Statement}

\begin{theorem}[SU(2) Mass-Gap Theorem]
\label{thm:mass-gap-su2}
Let $\Hop$ be the self-adjoint Hamiltonian on the physical Hilbert
space $\Hphys$ constructed by the Osterwalder--Schrader reconstruction
of the continuum SU(2) Wilson measure $\muinf$
(\Cref{ch:os-reconstruction} Theorem~6.2).  Then there
exists an explicit, computable, $Q'$-rational $m_0 > 0$ such that
\begin{equation*}
  \sigma(\Hop) \cap (0, m_0) \;=\; \emptyset.
\end{equation*}
Equivalently, the \emph{mass} $m := \inf \sigma_+(\Hop)$, where
$\sigma_+(\Hop) := \sigma(\Hop) \setminus \{0\}$, satisfies
\begin{equation*}
  m \;\ge\; m_0 \;>\; 0.
\end{equation*}
The constructed $m_0$ depends on the choice of heat-kernel convention;
canonical and Wilson conventions give different but both nonzero
explicit values per \S\ref{sec:explicit-bounds}.
\end{theorem}

\subsection{Comparison with the Clay statement}

The Clay Millennium Problem statement ([JaffeWitten2006],
\texttt{docs/pdf\_references/JaffeWitten2006.pdf}) reads:

\begin{quote}
\emph{Yang--Mills Existence and Mass Gap.} Prove that for any
compact simple gauge group $G$, a non-trivial quantum Yang--Mills
theory exists on $\mathbb{R}^4$ and has a mass gap $\Delta > 0$.
Existence includes establishing axiomatic properties at least as
strong as those cited in [Streater--Wightman, Osterwalder--Schrader].

\medskip

A quantum field theory has a mass gap $\Delta$ if $H$ has no
spectrum in the interval $(0, \Delta)$ for some $\Delta > 0$. The
supremum of such $\Delta$ is the mass $m$, and we require
$m < \infty$.
\end{quote}

Our Theorem~\ref{thm:mass-gap-su2} addresses the SU(2) instance,
with:

\begin{itemize}
  \item \textbf{Existence of the QFT}: provided by the projective-
    limit construction of $\muinf$ (\Cref{ch:continuum-measure}).
  \item \textbf{Axiomatic properties at least as strong as
    Osterwalder--Schrader}: provided by the OS-axiom verification of
    \Cref{ch:os-axioms}.  The full
    Streater--Wightman field reconstruction is a further analytical
    step not addressed here; we deliver the OS framework, which is
    what the Clay-statement axiomatic threshold requires.
  \item \textbf{Mass gap $\Delta > 0$}: provided by
    Theorem~\ref{thm:mass-gap-su2} with the explicit lower bound
    $m_0 > 0$.  Our $m_0$ is a constructive Q'-rational lower bound;
    finiteness of $m$ ($m < \infty$) is automatic from the
    continuum-limit construction (the Hamiltonian on the
    reconstructed Hilbert space has bounded spectrum on every
    finite-energy cylinder).
\end{itemize}

\textbf{Generality of the constructed $m_0$.}  The lower bound $m_0$
is constructive: it is computed from the Phase-3 pilot numerical
data via the closed chain $\sighinf \ge \varepsilon \Rightarrow m
\ge m_0$ derived below.  This is stronger than a pure-existence
mass-gap claim; it produces an explicit $Q'$-rational number with
the property that any spectral value of $\Hop$ in $(0, m_0)$ is
impossible.

\section{The proof: composition of Milestones 7.1, 8.1, 8.2}
\label{sec:proof}

The proof is a three-step composition:

\begin{enumerate}
  \item \textbf{the $\sighinf$-positivity analysis input}: the
    renormalized string tension is bounded below,
    $\sighinf \ge \varepsilon > 0$ in $Q'$.
  \item \textbf{the OS-axiom verification bridge}: $\sighinf > 0$
    gives the Wilson-loop area law for $\muinf$ and the OS4 cluster
    property with a strictly positive mass rate $m > 0$.
  \item \textbf{the OS-reconstruction analysis bridge}: the
    area-law--to--gap bridge plus the cluster-to-gap theorem give the
    spectral gap $m \ge \sqrt{\sighinf}$ (in $\sigma_0 = 1$ units).
    The square root is the dimensional content of the
    string-tension--to--mass-gap relation: $\sighinf$ is a string
    tension (inverse area), $m$ a mass (inverse length).
\end{enumerate}

\subsection{Step 1: $\sighinf \ge \varepsilon > 0$ (the $\sighinf$-positivity analysis)}

\begin{lemma}[the $\sighinf$-positivity analysis framework theorem, in summary]
\label{lem:sigh-positive}
There exists a $Q'$-rational $\varepsilon > 0$ such that for every
$N \in \mathbb{N}$, the running renormalized string tension
$\sigh_N := \sigh_0 - \sum_{k < N} \eta_k$ along the UV-direction
$\sigh$-trajectory satisfies $\sigh_N \ge \varepsilon$.
Equivalently, $\sigh_0 - \mathrm{bound} \ge \varepsilon$ where
$\mathrm{bound}$ is any uniform partial-sum upper bound on the
$\eta$-sequence.
\end{lemma}

\begin{proof}
This is \texttt{lean/YMGap/Continuum/SigmaHatLimit.lean}
\texttt{sigma\_hat\_uv\_limit\_positive}, with the framework theorem
$\texttt{sigma\_hat\_uniform\_positive\_lower\_bound}$ composing the
uniform $\eta$-bound (from any of the Phase-6
$\sigma$-budget-closure providers) with the gap condition
$\mathrm{bound} + \varepsilon \le \sigh_0$.
\end{proof}

The constructive content of Lemma~\ref{lem:sigh-positive}: from
the Phase-3 pilot rational data
(\codeanchor{computations/q\_lift.py::pilot\_eta\_rational\_bound}),
the gap $\varepsilon$ is exactly
$\sigh_0 - \mathrm{eta\_total\_pilot}$ in the post-$\clambda$
Casimir-gap upgrade:
\begin{align*}
  \text{(canonical, } g_0 = 0.30\text{):} \qquad
    \varepsilon &= 1 - 9/10000 = 9991/10000, \\
  \text{(Wilson, } g_0 = 0.50\text{):} \qquad
    \varepsilon &= 1 - 3/10000 = 9997/10000.
\end{align*}
Both are strictly positive $\Qp$-rationals; the canonical case is
slightly worse (more $\eta$-budget consumed) because the canonical
pilot operates closer to the small-field convergence threshold.
The previous values $9967/10000$ (canonical) and $9988/10000$
(Wilson) were tracked when $\clambda$ was a fiducial $\tfrac12$;
the upgrade to the derived Casimir gap $\clambda = \tfrac34$
(\Cref{ch:small-field-kp},
Theorem~\ref{thm:c-lambda-casimir}) tightens the budget by a
factor of $\approx 3$ at both operating points.

\subsection{Step 2: $\sighinf > 0$ $\Rightarrow$ OS4 cluster (the OS-axiom verification)}

\begin{lemma}[OS4 cluster with a positive mass rate]
\label{lem:os4-cluster}
There is a mass rate $m > 0$ such that for test functions $F, G$ with
spatial supports separated by Euclidean distance $d \ge a_0$,
\[
  \bigl| \langle F G \rangle_{\muinf} - \langle F \rangle_{\muinf}
                                          \langle G \rangle_{\muinf} \bigr|
  \;\le\;
  C_{F, G} \, \exp(-m \cdot d).
\]
The positivity $m > 0$ follows from $\sighinf \ge \varepsilon > 0$
(Lemma~\ref{lem:sigh-positive}).
\end{lemma}

\begin{proof}
\Cref{ch:os-axioms} \cref{thm:os4} (OS4 cluster property), via the
convergent polymer-cluster expansion of \Cref{ch:sheaf-framework}
\S7.  The cluster rate is a mass (inverse length); its
string-tension-derived lower bound is deferred to
Lemma~\ref{lem:cluster-to-gap}.
\end{proof}

\subsection{Step 3: cluster $\Rightarrow$ spectral gap (the OS-reconstruction analysis)}

\begin{lemma}[Area-law--to--gap bridge $+$ cluster-to-gap]
\label{lem:cluster-to-gap}
The OS Hamiltonian $\Hop$ on $\Hphys$ satisfies
\[
  \sigma(\Hop) \cap \bigl(0, \sqrt{\sighinf}\,\bigr) \;=\; \emptyset,
  \qquad \text{i.e., } \;\; m \;\ge\; \sqrt{\sighinf}
  \quad (\sigma_0 = 1 \text{ units}).
\]
\end{lemma}

\begin{proof}
\Cref{ch:os-reconstruction} \cref{cor:gap-from-sigma}: the
cluster-to-gap theorem (\cref{thm:cluster-to-gap}) gives
$m \ge m_{\mathrm c}$ for the cluster rate $m_{\mathrm c}$ of
Lemma~\ref{lem:os4-cluster} via the Källén--Lehmann representation of
the time-direction two-point function, and the area-law--to--gap
bridge (\cref{thm:area-law-to-gap-bridge}) gives
$m_{\mathrm c} \ge \sqrt{\sighinf}$.  The square root is the
dimensional content of the relation (string tension $\sighinf$ is
inverse area, the gap $m$ inverse length); the $O(1)$ constant is
taken conservatively as $1$, the looseness analysed in
\S\ref{sec:sharpness}.
\end{proof}

\subsection{Composition: proof of Theorem~\ref{thm:mass-gap-su2}}

\begin{proof}[Proof of Theorem~\ref{thm:mass-gap-su2}]
By Lemma~\ref{lem:sigh-positive}, the renormalized string tension
satisfies $\sighinf \ge \varepsilon > 0$ for $\varepsilon$ given
explicitly by the Phase-3 pilot rational data.  By
Lemma~\ref{lem:os4-cluster}, $\muinf$ satisfies OS4 with a strictly
positive mass rate.  By Lemma~\ref{lem:cluster-to-gap} (the
area-law--to--gap bridge composed with cluster-to-gap), the OS
Hamiltonian satisfies $\sigma(\Hop) \cap (0, \sqrt{\sighinf}) =
\emptyset$, hence
\begin{equation*}
  m \;\ge\; \sqrt{\sighinf} \;\ge\; \sqrt{\varepsilon} \;>\; 0
  \qquad (\sigma_0 = 1 \text{ units}).
\end{equation*}
Accordingly define the constructive lower bound
$m_0 := \lfloor \sqrt{\varepsilon} \rfloor_{Q'}$, the largest
$Q'$-rational with $m_0^2 \le \varepsilon$; then $m_0 \le
\sqrt{\varepsilon} \le m$ and $m_0 > 0$.  The square root is the
dimensional content of the string-tension--to--mass-gap relation
(Lemma~\ref{lem:cluster-to-gap}), not a change of units; the gap
between the conservative unit constant and the physical ratio
$m/\sqrt{\sigma}\approx 3.74$ is analysed in \S\ref{sec:sharpness}.

The explicit $m_0$ values are computed in \S\ref{sec:explicit-bounds}.
\end{proof}

\section{Explicit constructive bounds $m_0$}
\label{sec:explicit-bounds}

The proof of Theorem~\ref{thm:mass-gap-su2} reduces the mass-gap
lower bound to extracting a $Q'$-rational $m_0$ with
$m_0^2 \le \sighinf$ (so that $m_0 \le \sqrt{\sighinf} \le m$ by the
area-law--to--gap bridge), where $\sighinf$ is the Phase-7.1
renormalized-string-tension lower bound.

\begin{proposition}[Explicit $\massgap$ from the canonical pilot]
\label{prop:m0-canonical}
For the canonical-convention $\SUtwo$ pilot at $g_0 = 0.30$,
$K = 3$ transitions, denominator $10000$, the mass-gap lower bound is
\begin{equation*}
  \massgap^{(\mathrm{can})} \;:=\; \frac{999}{1000} \;\in\; \Qp,
  \qquad
  \massgap^{(\mathrm{can})} \;\le\; m.
\end{equation*}
\end{proposition}

\begin{proof}
We verify $\bigl(999/1000\bigr)^2 \le \varepsilon$:
\[
  \bigl(999/1000\bigr)^2 \;=\; \frac{998001}{1000000} \;=\; 0.998001,
\]
which is less than $\varepsilon = 9991/10000 = 0.9991$.  The
inequality $0.998001 \le 0.9991$ is decidable by
\texttt{decide} on the cross-product representation
$(999 \cdot 999 \cdot 10000 \le 9991 \cdot 1000 \cdot 1000)$,
i.e., $9980010000 \le 9991000000$.

By Theorem~\ref{thm:mass-gap-su2}'s constructive content,
$\massgap^{(\mathrm{can})} \le m$.
\end{proof}

\begin{proposition}[Explicit $\massgap$ from the Wilson pilot]
\label{prop:m0-wilson}
For the Wilson-convention $\SUtwo$ pilot at $g_0 = 0.50$, $K = 3$
transitions, denominator $10000$, the mass-gap lower bound is
\begin{equation*}
  \massgap^{(\mathrm{wil})} \;:=\; \frac{999}{1000} \;\in\; \Qp,
  \qquad
  \massgap^{(\mathrm{wil})} \;\le\; m.
\end{equation*}
\end{proposition}

\begin{proof}
$\bigl(999/1000\bigr)^2 = 998001/1000000 = 0.998001 \le \varepsilon =
9997/10000 = 0.9997$ by
$998001 \cdot 10000 \le 9997 \cdot 1000000$, i.e.,
$9980010000 \le 9997000000$.  Decidable on the cross-product
representation.
\end{proof}

\subsection{Units of $m_0$ and the meaning of "natural units"}

The $m_0$ values in Propositions \ref{prop:m0-canonical}--\ref{prop:m0-wilson}
are dimensionless: they live in $Q'$.  Two ways to read them:

\begin{enumerate}
  \item \textbf{Units of $\sigh_0^{1/2}$}: with $\sigh_0 = 1$ chosen
    as the unit at scale $0$, $m_0 \in (0, 1)$ is a fraction of the
    bare-lattice string-tension scale.  The bounds
    $\massgap^{(\mathrm{can})} = 0.999$ and $\massgap^{(\mathrm{wil})} = 0.999$
    are dimensionless and read directly.
  \item \textbf{Physical units}: pick a physical scale (e.g.\
    $\sqrt{\sigma_0} \approx 440\,\text{MeV}$ from the standard QCD
    string-tension calibration of [JaffeWitten2006] \S2; or
    equivalently $\sqrt{\sigma} \approx 400\,\text{MeV}$ under the
    pure-gauge calibration common in the lattice literature).
    Then $m_0^{(\mathrm{can})} \cdot \sqrt{\sigma_0} \approx 0.999
    \cdot 440\,\text{MeV} \approx 440\,\text{MeV}$ in the canonical
    case.  In the SU(2)-Yang--Mills setting where there is no
    analog of QCD's pion (the physical mass is the lightest
    glueball), this is a glueball-mass lower bound of order
    $440\,\text{MeV}$; see \S\ref{sec:sharpness-teper} for the
    explicit comparison to Teper's lattice spectroscopy.
\end{enumerate}

The dimensionless form is the constructive content of this
milestone; the physical-unit translation is a calibration choice
made downstream.

\section{Sharpness and what remains}
\label{sec:sharpness}

\subsection{The proven $m_0$ is not sharp}

The bounds $m_0^{(\mathrm{can})} = 0.999$ and
$m_0^{(\mathrm{wil})} = 0.999$ are constructive but not sharp.  The
\emph{actual} SU(2) mass (the lightest excited-state energy of
$\Hop$) is conjecturally of order $\sqrt{\sigma_0} \cdot O(1)$, with
$O(1)$ coefficient depending on the glueball spectrum.  The proven
lower bound is $\sim 1.0$ in units of $\sqrt{\sigma_0}$, which is
roughly the right order of magnitude but quantitatively loose by
the $O(1)$ factor controlled by the leading $0^{++}$ glueball mass
of pure SU(2) gauge theory.  The next subsection makes this
quantitative against the lattice-spectroscopy reference of
\cite{Teper1998Glueballs} and \cite{LuciniTeperWenger2004}.

\subsection{Comparison to Teper's SU(2) glueball lattice spectroscopy}
\label{sec:sharpness-teper}

The standard lattice-spectroscopy result for pure SU(2) Yang--Mills
\cite{Teper1998Glueballs,LuciniTeperWenger2004} is that the lightest
state in the spectrum is the $0^{++}$ glueball, with continuum-
extrapolated dimensionless mass
\begin{equation}
  \label{eq:teper-su2}
  \frac{m(0^{++})}{\sqrt{\sigma}} \;\approx\; 3.74,
\end{equation}
where $\sigma$ is the SU(2) string tension and the value is from
\cite{Teper1998Glueballs} (with the improved-operator continuum
extrapolation of \cite{LuciniTeperWenger2004} yielding the consistent
$3.78 \pm 0.07$).  Translating to physical units via the standard
pure-gauge calibration $\sqrt{\sigma} \approx 400$--$440\,\text{MeV}$:
\begin{equation*}
  m(0^{++})_{\text{SU(2)}} \;\approx\; 1.5\,\text{GeV}
  \quad\text{to}\quad 1.65\,\text{GeV}.
\end{equation*}
The constructive lower bound of Theorem~\ref{thm:mass-gap-su2} (in
the same physical units) is
\begin{equation*}
  m_0^{\text{phys}}
  \;=\; m_0^{(\mathrm{can})} \cdot \sqrt{\sigma_0}
  \;\approx\; 0.999 \cdot 440\,\text{MeV}
  \;\approx\; 440\,\text{MeV}.
\end{equation*}

\paragraph{The ratio.} The proven lower bound is loose relative to
the conjectured actual mass by the factor
\begin{equation}
  \label{eq:teper-ratio}
  \frac{m(0^{++})_{\text{SU(2)}}}{m_0^{\text{phys}}}
  \;\approx\; \frac{1500\,\text{MeV}}{440\,\text{MeV}}
  \;\approx\; 3.4
\end{equation}
(under the $\sqrt{\sigma} \approx 400\,\text{MeV}$ calibration; the
ratio rises to $\sim 3.75$ under the $\sqrt{\sigma} \approx 440\,
\text{MeV}$ calibration of \cite{JaffeWitten2006}).  Either
calibration yields a ratio in the range $[3.4, 3.8]$.

\paragraph{The interpretation: the lower bound is meaningful within
an $O(1)$ factor.}  Equation \eqref{eq:teper-ratio} says the
constructive $m_0$ is loose by a factor $\sim 3.4$ relative to the
conjectured actual mass.  Crucially, the bound is \emph{within an
$O(1)$ factor} of the truth — not off by orders of magnitude.  Three
qualifying observations:

\begin{enumerate}
  \item \textbf{The bound is in the right physical regime.}  The
    Clay Millennium problem asks for $\Delta > 0$ \emph{without
    quantitative constraint}.  Theorem~\ref{thm:mass-gap-su2}
    delivers $\Delta \ge m_0 \approx 440\,\text{MeV}$, which is
    $44$\% of the conjectured actual mass and $0.999$ in
    dimensionless units of $\sqrt{\sigma_0}$.  This is the right
    physical scale (hadronic-mass regime, not Planck-scale or
    QED-scale), the result of the $\sqrt{\sigma_0}$ identification
    in the OS-reconstruction $\sigh \to m$ step
    (\S\ref{sec:proof}).
  \item \textbf{The factor $3.4$ is not unbounded.}  The next-
    order-of-magnitude lattice-physics-side errors (statistical:
    $\pm 0.07$ from \cite{LuciniTeperWenger2004}; systematic:
    $\sim 10$\% from continuum extrapolation; calibration:
    $\pm 40\,\text{MeV}$ from string-tension uncertainty in
    SU(2)-vs-QCD comparison) collectively shift the upper-side
    benchmark by less than $\pm 20$\%.  The $\sim 3.4$ factor is
    therefore robust within standard lattice systematics.
  \item \textbf{The recommended refinement priorities below address
    the $\sim 3.4$ factor.}  The Brydges--Federbush
    cluster-decay-to-spectral-gap step (the third item of the
    "Recommended refinement priorities" subsection below) could
    tighten the $m \ge \sighinf^{1/2}$ identification by an $O(1)$
    multiplicative factor; combined with the Phase-6.5
    chaining-tightening propagated into $\eta$ (first item), the
    loose factor could plausibly be reduced to the $1.5$--$2$ range,
    closer to the conjectured sharp value.
\end{enumerate}

The standard reading: \emph{Theorem~\ref{thm:mass-gap-su2} establishes
the constructive existence of a strictly positive mass gap whose
quantitative lower bound is meaningful within an $O(1)$ factor of
the conjectured truth, with the loose factor tied to identifiable
analytical refinements outside the existence-proof's load-bearing
chain.}  This is the strongest claim the proof supports; it is
neither a sharp-value claim nor a vacuous-existence claim.

\paragraph{Why a non-sharp bound is still meaningful.}  Three points
make the looseness a feature of the \emph{calibration}, not a weakness
of the \emph{result}.
\begin{enumerate}[label=(\alph*)]
  \item \textbf{The Clay condition is $\Delta > 0$, not a value.}  The
    Millennium Problem asks for a mass gap $\Delta > 0$ with no
    quantitative target; $\massgap \ge 999/1000$ meets that condition
    and over-delivers by naming an explicit positive constant.
  \item \textbf{It is, to our knowledge, the first \emph{explicit}
    constructive lower bound.}  Prior constructive control of
    four-dimensional gauge theory (Bałaban; Magnen--Rivasseau--Sénéor)
    establishes that a positive gap \emph{exists} in suitable regimes
    but does not collapse the renormalized string tension to a named
    number; here every constant in the chain
    $\sighinf \ge \varepsilon \Rightarrow m \ge \sqrt{\varepsilon}$ is
    an explicit $\Qp$-rational, and the final inequality is
    machine-checked (\Cref{ch:mass-gap-theorem} Lean witness,
    axiom-gate $\codepath{[propext]}$).
  \item \textbf{The gap is an $O(1)$ factor from the truth, not orders
    of magnitude.}  The loose factor is the single
    string-tension--to--mass-gap constant ($\approx 3.74$,
    \cref{sec:looseness-decomposition}); the bound lands in the
    correct hadronic-mass regime ($\approx 440\,\text{MeV}$), within
    a factor $\sim 3.7$ of the conjectured glueball mass --- the
    difference between a meaningful constructive bound and a vacuous
    one.
\end{enumerate}

\paragraph{Side-note on the calibration.}  The two conventional
calibrations $\sqrt{\sigma} \approx 400$ and $440\,\text{MeV}$
differ by $\sim 10$\% reflecting the SU(2)-vs-QCD scale-setting
choice: the higher value uses the experimental hadronic-mass
calibration of QCD (where $\sqrt{\sigma}$ is determined from
quark--antiquark linear potential fits), while the lower value uses
the pure-gauge SU(2) calibration that \cite{Teper1998Glueballs}
reports.  This $10$\% spread is irrelevant to the $O(1)$ framing
above.

\subsection{Sources of looseness in the lower bound}
\label{sec:looseness-decomposition}

The $\sim 3.74$ gap between $m_0$ and the conjectured $m$ factorizes
across the proof chain.  The decisive observation is that it is
\emph{dominated by a single step} --- the area-law--to--mass-gap
bridge --- while every $\sigh$-trajectory-side source contributes
below the $2\%$ level, because $\sighinf \approx 0.999$ already leaves
almost no headroom and $m_0 = \sqrt{\sighinf}$ compresses it further.

\begin{center}
\renewcommand{\arraystretch}{1.3}
\begin{tabular}{p{3.5cm}p{2.6cm}p{2.3cm}p{3.0cm}}
\toprule
\textbf{Source} & \textbf{Enters via} & \textbf{Magnitude} &
  \textbf{Refinement status} \\
\midrule
Linear-decay-to-gap ($m$ vs $\sqrt{\sigma}$) &
  cluster-to-gap bridge, \cref{thm:area-law-to-gap-bridge} &
  $\sim 3.74\times$ \textbf{(dominant)} &
  open: sharper $O(1)$ constant, Brydges--Federbush 1981 \\
$\Kcap$ union bound ($\Kcap \le \Ksf + \Klf$) &
  overlap KP, \cref{ch:overlap-kp} &
  $1.01$--$1.18\times$ on $\Kcap$ &
  empirically tight (Conj.\ 7.4, $1.014$ at $g_0=0.30$); $K$-level
  Lean-closed \\
$\cpolymer$ overcounting ($2$ vs $1$) &
  $\eta = (\cpolymer/\Amin)\Dcap$ &
  $\le 2\times$ on $\eta$;\, $\sim\!+5\times10^{-4}$ on $\sighinf$ &
  open: sign-preservation lemma \\
$\Dcap$ buffer growth &
  same $\eta$ coupling &
  true UV value $0$;\, surrogate $\le 1.2\times10^{-2}$ &
  structural: $\eta \equiv 0$ in the UV (\cref{rmk:explicit-tail}) \\
Polymer-size cutoff / chaining &
  KP sum truncation; $c_{\mathrm{chain}}(g)$ &
  $c_{\mathrm{chain}}\!\in\!(0,1]$ \emph{uplift};\, $\sim\!+10^{-3}$ on $m_0$ &
  tightening, applied (\cref{ch:chaining-estimate}) \\
\bottomrule
\end{tabular}
\end{center}

Reading the table: the four $\sigh$-side rows compound to under
$1.02\times$ --- they move $\sighinf$ only in its third--fourth
decimal, hence $m_0$ in its fourth --- so \emph{the entire $\sim 3.74$
looseness is the linear-decay-to-gap constant}.  Two consequences:

\begin{enumerate}
  \item \textbf{The only lever that matters is the $\sqrt{\sigma}$
    bridge.}  Tightening the union bound, $\cpolymer$, $\Dcap$, or the
    polymer cutoff cannot close the gap: their combined headroom is
    $<2\%$.  A sharper area-law--to--gap relation
    (\cref{thm:area-law-to-gap-bridge}, conservative constant $1$ vs
    the physical $\approx 3.74$) is the sole route to a quantitatively
    sharp bound.
  \item \textbf{Pilot truncation is not a looseness source.}  Running
    the pilot to higher $K$ \emph{consumes} more surrogate $\eta$ and
    would \emph{decrease} $\varepsilon$; but the true UV-direction
    trajectory has $\eta \equiv 0$ (\cref{rmk:explicit-tail}), so
    $\sighinf = \sigma_0 = 1$ and the finite-$K=3$ quotable
    $9991/10000$ is already within $9\times10^{-4}$ of the
    method's ceiling $m_0 \le 1$.
\end{enumerate}

\subsection{Recommended refinement priorities}

In decreasing order of impact:

\begin{itemize}
  \item Propagate Phase-6.5 chaining tightening into the $\eta$
    derivation of \texttt{computations/sigma\_hat\_trajectory.py}.
    Estimated $m_0$ uplift: factor of $\sim 1.001$ canonical,
    $\sim 1.0005$ Wilson (modest, because $\mathrm{eta\_total}$ is
    already small in absolute terms).
  \item Refine the $c_\lambda$ fiducial of the small-field
    certificate (flagged refinement target in the fiducial audit),
    which would tighten the per-step $\eta$ at the upstream level.
    Estimated $m_0$ uplift: indirect, dependent on the refinement
    magnitude.
  \item Replace the $m \ge \sighinf^{1/2}$ linear-decay-to-gap step
    with a sharper Brydges--Federbush-style estimate.  Estimated
    $m_0$ uplift: linear in the relation, $\sim 1.0$--$2.0$ range.
\end{itemize}

These are all open analytical refinements, NOT blockers for the
existence statement of Theorem~\ref{thm:mass-gap-su2}.

\section{What this milestone does not deliver}
\label{sec:scope}

\begin{itemize}
  \item \textbf{Sharp mass-gap value.}  The bound $\massgap \ge 0.999$
    canonical is constructive but not sharp; the actual $m$ is
    larger.  Refining the bound is part of the
    $\SUN$-extension and analytic-refinement programme.
  \item \textbf{Lean formalization of the statement.}  The Lean
    module \codepath{lean/YMGap/MassGap/SU2.lean} carries the
    Lean-formal statement of the mass-gap theorem, including the
    chain from \codepath{sigma\_hat\_uv\_limit\_positive} to the
    spectral statement.  This chapter supplies the paper-level
    proof; the Lean module formalizes the statement against the
    abstract type signatures of the continuum-limit and OS-axiom
    layers.
  \item \textbf{SU(3) extension.}  the SU(N) extension framework.
  \item \textbf{Glueball-spectrum identification.}  Beyond the
    Clay-statement requirement; future-work.
\end{itemize}

\section*{References}

\noindent[Bal87] T.~Bałaban, ``Renormalization group approach to
lattice gauge field theories.~I,'' Comm.\ Math.\ Phys.\ \textbf{109}
(1987), 249--301.

\noindent[BF81] D.~Brydges, P.~Federbush, ``A new form of the
Mayer expansion in classical statistical mechanics,'' J.\ Math.\ Phys.\
\textbf{22} (1981), 2064--2074.  Sharper mass-gap-from-clustering
relations.

\noindent[GJ87] J.~Glimm, A.~Jaffe, \emph{Quantum Physics: A
Functional Integral Point of View}, 2nd ed., Springer 1987.  \S6.1
OS axioms + reconstruction; \S20 polymer cluster expansions.

\noindent[JW06] A.~Jaffe, E.~Witten, ``Quantum Yang--Mills theory,''
Clay Millennium Problem statement (2006).

\noindent[OS73] K.~Osterwalder, R.~Schrader, ``Axioms for Euclidean
Green's functions,'' Comm.\ Math.\ Phys.\ \textbf{31} (1973), 83--112.

\noindent[OS75] K.~Osterwalder, R.~Schrader, ``Axioms for Euclidean
Green's functions.~II,'' Comm.\ Math.\ Phys.\ \textbf{42} (1975),
281--305.

\noindent[Sim74] B.~Simon, \emph{The $P(\phi)_2$ Euclidean (Quantum)
Field Theory}, Princeton 1974.

\noindent[Tep98] M.~Teper, ``Glueball masses and other physical
properties of SU(N) gauge theories in $d=3+1$: a review of lattice
results for theorists,'' arXiv:hep-th/9812187 (1998).  SU(2) review
value: $m(0^{++}) / \sqrt{\sigma} \approx 3.74$.

\noindent[LTW04] B.~Lucini, M.~Teper, U.~Wenger, ``Glueballs and
$k$-strings in SU(N) gauge theories: calculations with improved
operators,'' JHEP \textbf{0406} (2004) 012.  Improved-operator SU(2)
continuum extrapolation: $m(0^{++}) / \sqrt{\sigma} \approx 3.78
\pm 0.07$.
